% Section 13 --- The two processions: 2 Kings 6 with Paralipomenon, the Uzzah
% crux, Obededom, the dance, and the synoptic table. Drafted per
% DRAFTING-BRIEF.md sec. 9 from ot-corpus.md secs. 30-33, 39-44, VIII;
% w2-textcrit.md secs. 1, 3; w2-rights.md secs. 1.14-1.17, 1.19-1.21, 2.
\sectionguard
\section{The Two Processions}

On the road to Jerusalem a man put out his hand to steady the ark of God, and
God struck him dead. Every honest reader stumbles at that sentence, and ought
to. Oza's motive, so far as the text lets us see one at all, was to help: the
oxen kicked, the ark leaned, and he took hold of it. The instinct that
protests on his behalf is not irreligion; it is a sense of justice, and it
deserves to be stated at full strength before anything else in this chapter is
said. A God who strikes dead the man who steadies the leaning sacred object
looks, on a first reading, like caprice.

But the chapter of Scripture in which Oza dies does not contain one
procession. It contains two. Three months apart, over the same last miles,
with the same ark and the same king, Israel makes the same journey twice ---
and the first ends in death and the second in dancing. Whatever changed
between them is the thing the narrative wants found, and it takes some care
to find it, because the story has come down to us in two tellings: the
historian's account in 2 Kings 6 (= 2 Samuel 6, in the Hebrew naming), and the
liturgist's account in 1 Paralipomenon 13 and 15--16 (= 1 Chronicles). This
study prints both and harmonizes neither. Where they differ, the differences
are themselves evidence.

\subsection{The first procession: the new cart}

Samuel gives no motive. After the long years at Cariathiarim, with ``all the
chosen men of Israel, thirty thousand,'' David simply acts:

\begin{witnessbox}{2 Kings 6:2, Douay-Rheims}
And David arose and went, with all the people that were with him of the men
of Juda to fetch the ark of God, upon which the name of the Lord of hosts is
invoked, who sitteth over it upon the cherubims.
\end{witnessbox}

It is as full a title as the ark bears anywhere in the histories: the Name
invoked over it, the throne above it. The narrator states what is being
moved in the same breath as the moving of it --- and the reader who has come
through Sinai and Bethsames already knows what those titles cost.\endnote{2
Kings 6, Douay-Rheims (Challoner), \url{https://drbo.org/chapter/10006.htm}.
Douay 6:2 reads ``of the men of Juda'' (Vulgate \term{de viris Juda}) where
the received Hebrew has what modern versions print as the place-name
Baale-judah; the oldest Samuel manuscript, 4QSam-a, reads a place-name here
outright --- Baalah, there identified with Cariathiarim (the identifying
phrase partly restored in the scroll), as at 1 Paralipomenon 13:6 (\work{The Dead Sea Scrolls Bible}, ed.\ Abegg, Flint,
and Ulrich, 1999, notes to 2 Samuel 6, read at archive.org).}
The Chronicler, writing later and for the temple's people, supplies
the motive Samuel withheld --- ``let us bring again the ark of our God to us:
for we sought it not in the days of Saul'' --- and widens the muster from an
army to a nation, ``from Sihor of Egypt, even to the entering into
Emath.''\endnote{1 Paralipomenon 13:3, 5, Douay-Rheims,
\url{https://drbo.org/chapter/13013.htm}.}

Then the method. They set the ark of God upon a new cart and brought it
``out of the house of Abinadab, who was in Gabaa,'' with Oza and Ahio to
drive it. A new cart drawn by oxen is exactly how the ark
had last travelled --- it is the transport the Philistine diviners devised at
the end of the seven months, a contrivance for a god whose people were not
there to carry him. Israel had another law. The sons of Caath were to bear
the covered ark on shoulders, and ``they shall not touch the vessels of the
sanctuary, lest they die.''\endnote{Numbers 4:15, Douay-Rheims,
\url{https://drbo.org/chapter/04004.htm}: ``then shall the sons of Caath
enter in to carry the things wrapped up: and they shall not touch the
vessels of the sanctuary, lest they die.'' The transport law is set out in
the Sinai chapter of this study; the Philistine cart, in the chapter on the
ark's captivity (1 Kings 6:7--9).} The cart was not a sin of malice; it was
a convenience. But a cart is how one moves furniture. Even we do not wheel
our dead to the grave while there are shoulders to bear them; some things
are carried because carrying is itself the confession of what they are. The
narrative does not say any of this at the outset. It lets the wheels turn.

At ``the floor of Nachon'' --- Paralipomenon calls it Chidon, and no two
witnesses quite agree on the name\endnote{The name of the threshing floor is
garbled in every direction: Nachon (Douay 2 Kings 6:6, with the received
Hebrew); Chidon (Douay 1 Paralipomenon 13:9); Nodab in the Vatican
manuscript of the Greek (Rahlfs, 2 Kingdoms 6:6, \term{heōs halō Nōdab},
read at blueletterbible.org); Nachor, with Achor as a variant, in the
hexaplaric witnesses; Aquila took it for no name at all but rendered ``the
prepared threshing floor'' (\term{halōnos hetoimēs}), a reading Field's note
finds preserved also in Eusebius's \work{Onomasticon}; the Samuel scroll
4QSam-a has yet another form, reported in the literature as Nodan; Josephus
has Cheidon (\work{Antiquities} 7.81, as reported). Field, \work{Origenis
Hexaplorum quae supersunt} (Oxford: Clarendon, 1875), vol.\ I, p.\ 555,
\url{https://archive.org/details/origenhexapla01unknuoft}; \work{The Dead
Sea Scrolls Bible}, notes to 2 Samuel 6. An unfamiliar name, variously
copied; no witness deserves the crown.} --- the oxen kicked, the ark leaned,
and ``Oza put forth his hand to the ark of God, and took hold of it.'' Then
the sentence on which the whole difficulty hangs: ``and he struck him for
his rashness: and he died there before the ark of God'' (2 Kings 6:6--7).

\subsection{The word that killed him}

``For his rashness'' reads like an explanation. It is important to know that
it is not one --- or rather, that it is the Douay's inheritance of an ancient
guess. At this point the received Hebrew has a phrase of famous obscurity,
\term{al-hashshal}, whose sense no one knows; English versions made on the
Hebrew have ventured ``for his error.'' The Vulgate rendered it \term{super
temeritate}, ``for the rashness,'' and the Douay follows. So the base text
of this study prints an interpretation of an obscure word, and the reader
should know he is reading one.

The ancient witnesses divide in an instructive way. The Old Greek of Samuel
had nothing at all where the phrase stands: in Field's reconstruction of the
Hexapla the column is empty, and the Greek phrase ``for the rashness''
(\term{epi tē propeteia}) enters the tradition only as a later addition
under Origen's asterisk, while Aquila rendered the Hebrew ``for the folly''
(\term{epi tē eknoia}).\endnote{Field, \work{Origenis Hexaplorum quae
supersunt}, vol.\ I, p.\ 555, at 2 Kings 6:7: the Old Greek is marked
\term{vacat}; the asterisked addition supplies \term{epi tē propeteia};
Aquila, \term{epi tē eknoia}.
\url{https://archive.org/details/origenhexapla01unknuoft}.} But the oldest
Hebrew manuscript of Samuel that survives, the Qumran scroll 4QSam-a ---
which is extant precisely across these verses --- reads a full clause where
the received text has the opaque word: God struck him for putting out his
hand to the ark, and he died there before God. That is, in substance, the
very text of Paralipomenon: ``struck him, because he had touched the ark; and
he died there before the Lord'' (1 Par 13:10). The common critical judgment
is that the received Hebrew of Samuel is the damaged remnant of that longer
clause; the honest counterweight is that this scroll often agrees with
Paralipomenon, so a harmonizing copyist cannot be excluded. Either way, the
evidence that the offense was the touch --- the hand on the holy thing ---
is ancient, Hebrew, and older than Christianity; it is not merely the
Chronicler's moral drawn after the fact.\endnote{\work{The Dead Sea Scrolls
Bible}, ed.\ Abegg, Flint, and Ulrich (1999), 2 Samuel 6 with its textual
notes (the scroll's readings at 6:6--7; the opening words of the 6:7 clause
are printed there restored in brackets), read at archive.org; the scroll
preserves 2 Samuel 6:2--9 and 6:12--18. The haplography judgment is
McCarter's (Anchor Bible 9, ad loc., as reported there); the letter-level
edition is DJD XVII (2005), not consulted for this study. Douay 1
Paralipomenon 13:9--10, \url{https://drbo.org/chapter/13013.htm}.}

And the touch had a law. Here the two tellings play their parts. Samuel
narrates and refuses to explain: the blow falls, David is angry, then
afraid, and the story moves on --- the older account tolerates the
inexplicable. The Chronicler, a century and more of reflection later, gives
the diagnosis: ``No one ought to carry the ark of God, but the Levites, whom
the Lord hath chosen to carry it, and to minister unto himself for ever'';
and, when the procession is finally done rightly, ``the sons of Levi took
the ark of God as Moses had commanded, according to the word of the Lord,
upon their shoulders, with the staves'' (1 Par 15:2, 15).\endnote{1
Paralipomenon 15, Douay-Rheims, \url{https://drbo.org/chapter/13015.htm}.
The Douay's 15:13 differs materially from versions made on the Hebrew: where
modern translations have David say the breach came ``because we did not
carry it the first time,'' the Vulgate and Douay read ``because you were not
present.'' The Chronicler's diagnosis therefore rests, in this study's base
text, on the verified wording of 15:2 and 15:15, not on 15:13.} The fault
was not a hand; it was a method. The cart put the ark where a stumble of
oxen could endanger it and a layman's reflex could touch it. Obedience would
have made the reflex impossible. What that severity means --- why the
Presence is not managed, and what kind of love the fear of the holy is ---
belongs to a later panel of this study. Here it is enough to see what the
texts themselves insist on: the second procession will change exactly one
thing, and it is the thing the law had commanded all along.

David's response ends where the men of Bethsames ended: in dread, and then
in distance --- only, in the king, an anger comes first that the narrative
records without explaining. ``How shall the ark of the Lord come
to me?'' (2 Kings 6:9). It is a question, not a refusal --- the one sentence
in the chapter that is still waiting for its full answer, and this study
will meet it again, on other lips, before the end.

\subsection{Three months in a borrowed house}

Unwilling to bring the ark home, David turns it aside: ``the ark of the Lord
abode in the house of Obededom the Gethite three months: and the Lord
blessed Obededom, and all his household'' (2 Kings 6:11). The detail is
quietly astonishing. The surname is, on its face, the name of Philistine
country: the ark
that had lately scourged the cities of its captors sits three months under
the roof of a man called after Geth, and nothing burns. The same Presence
that struck the floor of Nachon fills a private house with blessing. The
narrative offers no theory --- only the pairing, death on the road and
blessing in the house, set side by side and left there. The blessing is what
ends the suspension: David goes down to bring the ark up to his own city,
and this time everything is done differently.

\subsection{The second procession: shoulders and a dance}

This time the change of method arrives in a subordinate clause, as though it
were the most natural thing in the world: ``when they that carried the ark
of the Lord had gone six paces, he sacrificed an ox and a ram'' (2 Kings
6:13). Carried --- the verse's plainest word is its most important one. No
cart appears in the second procession; the ark goes up on men. The six paces
can be construed two ways: a single sacrifice after the first six steps, or
a sacrifice repeated at every sixth step, a road paved with offerings ---
and the ancient Greek versions render the verb in a tense that leaves the
repeated reading grammatically alive. Paralipomenon describes the
day independently: ``they offered in sacrifice seven oxen, and seven rams''
(1 Par 15:26). The two descriptions should not be merged, and there is a
manuscript to keep the editor humble: the same Qumran scroll of Samuel reads
seven bulls and seven rams in the Samuel verse itself.\endnote{Douay 2 Kings
6:13 (``an ox and a ram''); 1 Paralipomenon 15:26. The iterative
possibility: Aquila and Symmachus at 2 Kings 6:13 use the imperfect
\term{ethysiaze} (Field, vol.\ I, p.\ 555). 4QSam-a at 6:13 reads seven
bulls and seven rams, with Paralipomenon, against the received Hebrew and
the Greek (\work{The Dead Sea Scrolls Bible}, notes to 2 Samuel 6).}

Then the verse the whole tradition remembers: ``And David danced with all
his might before the Lord: and David was girded with a linen ephod'' (2
Kings 6:14). The king in a priest's plain linen, whirling before the ark ---
Paralipomenon, characteristically, robes the dancer: ``David was clothed
with a robe of fine linen,'' like the Levites around him, with the linen
ephod over it (1 Par 15:27). And at the window of the palace the procession
acquires its one cold spectator: ``Michol the daughter of Saul, looking out
through a window, saw king David leaping and dancing before the Lord: and
she despised him in her heart'' (2 Kings 6:16).\endnote{Symmachus, the
second-century Jewish translator, rendered David's whirling at this verse
\term{skirtōnta kai kanchazonta}, ``leaping and laughing'' --- \term{skirtan}
being the verb Saint Luke uses of the infant John in Elizabeth's womb. The
reading is preserved in two manuscripts (Codex Regius and the Chigi catena
manuscript, cod.\ 243): Field, \work{Origenis Hexaplorum quae supersunt},
vol.\ I, p.\ 555, note 28,
\url{https://archive.org/details/origenhexapla01unknuoft}. It is the one
ancient Greek Bible text in which David ``leaps'' before the ark --- and the
form is the participle \term{skirtōnta}, not the aorist sometimes quoted in
devotional use. Symmachus worked after Luke; what this parallel can and
cannot show is weighed in the panel on Saint Luke.} Samuel alone preserves
what follows: her derision of the king's self-abasement, his answer, and
then the terse notice that closes her story --- ``Therefore Michol the
daughter of Saul had no child to the day of her death'' (6:23). Paralipomenon
keeps the window and the scorn and drops everything after; the dialogue and
the childlessness are not in its account at all.

The day ends with the ark at rest: they ``set it in the midst of the tent,
which David had pitched for it,'' with holocausts and peace offerings; and
the Chronicler adds what for him is the whole point, a ministry that does
not disband when the procession does --- Asaph and his brethren left
``before the ark of the covenant of the Lord \ldots\ to minister in the
presence of the ark continually'' (1 Par 16:1, 37).\endnote{1 Paralipomenon
16, Douay-Rheims, \url{https://drbo.org/chapter/13016.htm}. In Paralipomenon
the procession is the founding act of Israel's ordered worship --- Levites,
singers, and gatekeepers organized around the ark (15:16--24), a psalm sung
that day, and a permanent ministry before the tent. That liturgical framing
is the Chronicler's own emphasis, and it is itself part of the ark's
reception history.} In Samuel a king brings a throne to his capital; in
Paralipomenon a liturgy is born.

\subsection{The two accounts, side by side}

The differences between the tellings are not blemishes to be sanded away;
they are two vantage points on one event, and the study's method requires
that both be printed as they stand. One caution about the table below: in
versions made from the received Hebrew, the phrase ``in her heart'' at the
window scene stands only in Paralipomenon, and commentators have made much
of the contrast; the Vulgate and Douay print the phrase in both verses, so
that contrast is invisible in this study's base text and is not tabled.

\begin{fourstable}{Point}{2 Kings 6}{1 Paralipomenon}{Note}
Motive & None stated; David simply acts (6:1--2). & ``For we sought it not
in the days of Saul'' (13:3). & The Chronicler passes the verdict on Saul
that Samuel withholds.\\
Assembly & ``The chosen men of Israel, thirty thousand'' (6:1). & All
Israel, ``from Sihor of Egypt, even to the entering into Emath'' (13:5). &
A picked army in Samuel; a whole nation at liturgy in Paralipomenon.\\
Threshing floor & Nachon (6:6). & Chidon (13:9). & An unfamiliar name,
copied differently in nearly every ancient witness; no reading deserves the
crown.\\
Why Oza died & ``For his rashness'' (6:7) --- the version's rendering of an
obscure Hebrew word. & ``Because he had touched the ark'' (13:10). & The
oldest Hebrew manuscript of Samuel reads with Paralipomenon here.\\
Sacrifice & ``An ox and a ram'' after six paces (6:13). & ``Seven oxen, and
seven rams'' (15:26). & Independent descriptions, not to be merged; a Qumran
Samuel scroll carries the Paralipomenon numbers.\\
David's dress & ``Girded with a linen ephod'' (6:14). & ``A robe of fine
linen,'' and the linen ephod (15:27). & The Chronicler robes the dancer.\\
Michol & Window scene, the exchange, the childlessness (6:16, 20--23). &
Window scene only (15:29). & Paralipomenon keeps her scorn and drops the
sequel.\\
Aftermath & Sacrifice, blessing, bread to all the people (6:17--19). &
Asaph and his brethren ``to minister in the presence of the ark
continually'' (16:37). & For the Chronicler the procession founds a
permanent liturgy.\\
\end{fourstable}

So the ark stands on Sion at last --- but in a tent. David himself is the
first to feel the unfinished edge of it: ``Dost thou see that I dwell in a
house of cedar, and the ark of God is lodged within skins?''\endnote{2 Kings
7:2, Douay-Rheims, \url{https://drbo.org/chapter/10007.htm} (Vulgate
\term{in medio pellium}, ``amid the skins,'' the tent-curtains); for
``Sion,'' 3 Kings 8:1, where the ark is carried ``out of the city of David,
that is, out of Sion'' (\url{https://drbo.org/chapter/11008.htm}).} The
answer he receives is one of the strangest courtesies in Scripture: God
declines the house, recalls that he has walked in a tabernacle and a tent
since Egypt, and promises that a son will build it. The king who danced must
leave the ark under canvas. A tent on Sion is not a resting place; it is a
pause. The house waits for the son.
